% Contenu partagé — dessins, tableaux, code source.

\chapter{Dessins, tableaux et code}%
\label{chap:graphics}
\minitoc%

%---------------------------------------------------------------------------
\section{Dessins}%
\label{sec:tikz}

Les aides TikZ du template posent une fenêtre, une grille et des axes sans
avoir à les redessiner à chaque figure.

\begin{center}
\begin{tikzpicture}[xmin=-2, xmax=2, ymin=-1.5, ymax=1.5, scale=1.2]
    \grille
    \axes{$t$}{$x$}
    \draw[thick, color=\ocotscolor{emph-a}, domain=-2:2, samples=100]
        plot (\x, {sin(\x r)});
\end{tikzpicture}
\end{center}

Les pointes de flèches par défaut viennent du template~:
\begin{center}
\begin{tikzpicture}
    \draw[->] (0,0) -- (3,0) node[right]{$x$};
    \draw[->] (0,-0.3) -- (0,1) node[above]{$y$};
\end{tikzpicture}
\end{center}

%---------------------------------------------------------------------------
\section{Annotation d'une image}%
\label{sec:tikzgraphics}

L'environnement \texttt{tikzgraphics} permet de placer des repères sur une
image bitmap en coordonnées pixel, sans recalcul d'échelle.

\begin{center}
\begin{tikzgraphics}{6cm}{2546}{552}{logo-insa}
    \pxnode[anchor=west, color=\ocotscolor{emph-b}]{2600}{276}{repere}{$\leftarrow$ repère}
\end{tikzgraphics}
\end{center}

%---------------------------------------------------------------------------
\section{Tableaux}%
\label{sec:tables}

Colonnes centrées et alignées à gauche de largeur fixe, filets d'épaisseurs
graduées, cales verticales.

\begin{center}
\begin{tabular}{C{3cm} C{3cm} C{3cm}}
    \bighrule
    Méthode & Ordre & Stabilité \Tstrut\Bstrut \\
    \medhrule
    Euler explicite  & 1 & conditionnelle \Tstrut \\
    Euler implicite  & 1 & inconditionnelle \\
    Runge--Kutta 4   & 4 & conditionnelle \Bstrut \\
    \bighrule
\end{tabular}
\end{center}

%---------------------------------------------------------------------------
\section{Code source}%
\label{sec:listings}

Le style de \texttt{listings} suit les couleurs du thème.

\begin{lstlisting}[language=Python, caption={Une méthode d'Euler explicite}]
def euler(f, t0, x0, tf, n):
    # pas de discretisation
    h = (tf - t0) / n
    t, x = t0, x0
    for k in range(n):
        x = x + h * f(t, x)
        t = t + h
    return t, x
\end{lstlisting}

%---------------------------------------------------------------------------
\section{Liens, URL, renvois}%
\label{sec:links}

Un lien nommé vers une
\href{https://github.com/ocourses/ocots-latex-template}{page du dépôt}, une URL
nue \url{https://ocourses.github.io/}, une note de bas de page qui renvoie
ailleurs\footnote{Voir aussi \url{https://www.ctan.org/pkg/tcolorbox}.}, et
\emphLink{un fragment coloré à la main} avec \verb|\emphLink|.

Les renvois internes traversent les chapitres~: le théorème~\ref{thm:cauchy} du
chapitre~\ref{chap:boxes}, la définition~\ref{def:ouvert}, la
proposition~\ref{prop:ademontrer} du chapitre~\ref{chap:blocks}, et
l'exercice~\ref{ex:matrices} du chapitre~\ref{chap:exercises}.

Index~: le mot \index{flot}flot est indexé, ainsi que le
mot \index{orbite}orbite.
